\documentclass[11pt,a4paper]{article}
\usepackage[margin=1in]{geometry}
\usepackage{amsmath,amssymb,amsthm,mathtools}
\usepackage[T1]{fontenc}
\usepackage{lmodern}
\usepackage{xcolor}
\usepackage[hidelinks]{hyperref}
\hypersetup{
  pdftitle={Exact Multiplicity Floors for Dual Singularities from Absorbing Gauss Fibres},
  pdfauthor={Lluis Eriksson}
}

\newtheorem{theorem}{Theorem}[section]
\newtheorem{proposition}[theorem]{Proposition}
\newtheorem{lemma}[theorem]{Lemma}
\newtheorem{corollary}[theorem]{Corollary}
\theoremstyle{definition}
\newtheorem{definition}[theorem]{Definition}
\newtheorem{remark}[theorem]{Remark}

\newcommand{\kk}{\mathbb C}
\newcommand{\PP}{\mathbb P}
\newcommand{\cO}{\mathcal O}
\newcommand{\cI}{\mathcal I}
\newcommand{\m}{\mathfrak m}
\newcommand{\Osc}{\operatorname{Osc}}
\newcommand{\mult}{\operatorname{mult}}
\newcommand{\Sing}{\operatorname{Sing}}
\newcommand{\Span}{\operatorname{Span}}

\title{Exact Multiplicity Floors for Dual Singularities\\
from Absorbing Gauss Fibres}
\author{Lluis Eriksson\\\small Independent researcher}
\date{August 24, 2026\quad (version 1.0)}

\begin{document}
\maketitle

\begin{abstract}
Let $X^d\subset\PP^{d+1}_{\mathbb C}$ be a smooth hypersurface, let
$\gamma_X:X\to X^\vee$ be its Gauss map, and let $Z$ be the reduced fibre
over a tangent hyperplane $W$.  In the complete $\cO_X(m)$ embedding, assume
that the span of $Z$ contains the order-$s$ osculating space at every support,
where $1\leq s\leq m$.  We prove the sharp dual-multiplicity floor
\[
 \mult_W(X^\vee)\geq s^d|Z|
 \geq s^d\binom{d+m}{d}.
\]
The first inequality is local and forced by absorption: the equation of $W$
restricts to order at least $s+1$ at every point of the fibre.  Each isolated
singularity of $X\cap W$ therefore has Milnor number at least $s^d$, while
the classical multiplicity--Milnor formula of Dimca, subsequently extended
by Parusi\'nski, identifies the multiplicity of the dual with the sum of
these Milnor numbers.  The number of analytic branches of
$X^\vee$ at $W$ is $|Z|$, so it has the same binomial lower floor.  When
$m\geq2s+1$, we also show that the rank-sensitive term from the predecessor
paper supplies no stricter inequality here, because a Gauss fibre lies in one
hyperplane.

The uniform bounds are attained simultaneously.  For every $(d,m,s)$ we construct a
smooth integral hypersurface with exactly
$\binom{d+m}{d}$ points in one Gauss fibre; their point span is proper and
order-$s$ absorbing, and the corresponding hyperplane section has an ordinary
isolated singularity of multiplicity $s+1$ at each point and no others.  The
dual point has exactly that many branches, affine tangent-cone cycle
$s^d\sum_{p\in Z}[T_{[W]}(p^\perp)]$, and multiplicity
$s^d\binom{d+m}{d}$.  For $s=1$ its tangent cone is a reduced arrangement of
distinct hyperplanes.  The construction is explicit only up to nonempty open
choices; no minimal hypersurface degree, analytic classification, or
exhaustive priority claim is made.
\end{abstract}

\section{Statement and scope}

Let $X^d\subset\PP^{d+1}_{\kk}$ be a smooth hypersurface of degree $D\geq2$,
let $H=\cO_X(1)$, and let
\[
 \gamma_X:X\longrightarrow (\PP^{d+1})^\vee
\]
be the Gauss morphism.  Its image is the dual hypersurface $X^\vee$.
For $\eta\in X^\vee$, write $W_\eta\subset\PP^{d+1}$ for the corresponding
hyperplane and
\[
 Z_\eta=\bigl(\gamma_X^{-1}(\eta)\bigr)_{\mathrm{red}}.
\]

Fix $m\geq1$.  For a nonempty finite reduced $Z\subset X$, let
\[
 S_Z=\operatorname{Im}\!\left(
 H^0(Z,H^m|_Z)^*\longrightarrow H^0(X,H^m)^*\right).
\]
For $p\in X$ and $a\geq0$, principal-parts evaluation gives the affine
osculating block
\[
 \widehat\Osc^a_p(H^m)=\operatorname{Im}\!\left(
 H^0((a+1)p,H^m|_{(a+1)p})^*\longrightarrow H^0(X,H^m)^*\right).
\]

\begin{definition}
For $1\leq s\leq m$, the set $Z$ is
\emph{point-span $s$-osculating-absorbing} if
\[
 \widehat\Osc^s_p(H^m)\subseteq S_Z\qquad(p\in Z).
\]
\end{definition}

Put
\[
 N_{d,m}=\binom{d+m}{d}.
\]

\begin{theorem}[Sharp dual-multiplicity floor]\label{thm:floor}
Let $\eta\in X^\vee$ and $Z=Z_\eta$.  If $Z$ is point-span
$s$-osculating-absorbing for the complete $H^m$ embedding, then:
\begin{enumerate}
\item $X^\vee$ has exactly $|Z|$ analytic branches at $\eta$ and
\[
 |Z|\geq N_{d,m};
\]
\item for every $p\in Z$, the germ of $X\cap W_\eta$ has multiplicity at
least $s+1$ and Milnor number $\mu_p\geq s^d$;
\item
\[
 \mult_\eta(X^\vee)=\sum_{p\in Z}\mu_p
 \geq s^d|Z|\geq s^dN_{d,m};
\]
\item as a cycle in $T_\eta((\PP^{d+1})^\vee)$, the affine tangent cone is
\[
 C_\eta(X^\vee)=\sum_{p\in Z}\mu_p[T_\eta(p^\perp)].
\]
\end{enumerate}
\end{theorem}

\begin{theorem}[Simultaneous sharpness]\label{thm:sharpness}
For every $d,m\geq1$, every $1\leq s\leq m$, and every
\[
 D\geq(s+2)N_{d,m}+1,
\]
there exist a smooth integral degree-$D$ hypersurface
$X^d\subset\PP^{d+1}$, a hyperplane $W$, and
$Z=(\gamma_X^{-1}([W]))_{\mathrm{red}}$ such that
\[
 |Z|=\dim S_Z=N_{d,m}<h^0(X,H^m),
\]
$Z$ is point-span $s$-osculating-absorbing, and $X\cap W$ has precisely the
points of $Z$ as singularities.  At every $p\in Z$ the section singularity is
ordinary of multiplicity $s+1$, so
\[
 \mu_p=s^d,\qquad
 \mult_{[W]}(X^\vee)=s^dN_{d,m},
\]
and
\[
 C_{[W]}(X^\vee)=s^d\sum_{p\in Z}[T_{[W]}(p^\perp)].
\]
In particular, every inequality in Theorem~\ref{thm:floor} is sharp.  When
$s=1$, the $N_{d,m}$ branches are smooth and
the tangent cone is the reduced union of the distinct hyperplanes $p^\perp$.
\end{theorem}

The multiplicity identity used here is classical
\cite{Dimca1986,Parusinski1991}; Dimca--Ilardi use the same bridge to study
nodal hyperplane sections and normal-crossing points of duals
\cite{DimcaIlardi2024}.  Accordingly, the factor $s^d$ is presented as an
explicit local consequence of that classical theory, not as a new general
principle about dual multiplicities.  The contribution isolated here is the
combination with point-span absorption: the resulting sharp binomial floor
and the exact dual geometry of equality
examples with proper point span.  This comparison is selective, not a
priority certification.

\section{Gauss fibres and branches}

\begin{lemma}[Normalization by the Gauss map]\label{lem:normalization}
The Gauss morphism $\gamma_X:X\to X^\vee$ is finite and birational.  Hence it
is the normalization of $X^\vee$, and the number of analytic branches of
$X^\vee$ at $\eta$ equals $|Z_\eta|$.
\end{lemma}

\begin{proof}
If $F$ defines $X$, its first partial derivatives have no common zero on $X$
because $X$ is smooth.  They define $\gamma_X$ and give
\[
 \gamma_X^*\cO_{X^\vee}(1)=\cO_X(D-1),
\]
which is ample.  A positive-dimensional fibre would make this line bundle
trivial on a positive-dimensional complete subvariety, contradicting
ampleness.  Thus $\gamma_X$ is finite.  Projective biduality identifies the
unique tangency point over a general smooth point of $X^\vee$, so the map is
birational.  Since $X$ is smooth and therefore normal, it is the
normalization.  Complex algebraic local rings are excellent, normalization
commutes with completion, and each completed local ring above $\eta$ is a
power-series domain.  The normalization fibre therefore indexes the analytic
branches.
\end{proof}

The equality
\[
 \Sing(X\cap W_\eta)=Z_\eta
\]
holds set-theoretically: a point of the hyperplane section is singular exactly
when $W_\eta$ is tangent to $X$ there.  Lemma~\ref{lem:normalization} makes
this set finite, so all section singularities are isolated.

\section{Absorption forces high contact}

\begin{lemma}[Local annihilator consequence]\label{lem:annihilator}
If $Z$ is point-span $s$-osculating-absorbing and
$a\in H^0(X,\cI_Z\otimes H^m)$, then
\[
 a_p\in\m_{X,p}^{s+1}H^m_p\qquad(p\in Z).
\]
\end{lemma}

\begin{proof}
The annihilator of $S_Z$ is $H^0(X,\cI_Z\otimes H^m)$.  The annihilator of
$\widehat\Osc^s_p(H^m)$ consists of sections whose germ vanishes modulo
$\m_{X,p}^{s+1}$.  The inclusion
$\widehat\Osc^s_p(H^m)\subseteq S_Z$ reverses under annihilators and gives the
claim.
\end{proof}

\begin{proposition}[Contact along an absorbing Gauss fibre]
\label{prop:contact}
Let $\ell_\eta$ be a linear equation for $W_\eta$.  Under the hypotheses of
Theorem~\ref{thm:floor},
\[
 (\ell_\eta|_X)_p\in\m_{X,p}^{s+1}H_p
 \qquad(p\in Z).
\]
\end{proposition}

\begin{proof}
Every point of $Z$ lies in $W_\eta$.  Fix $p\in Z$ and choose
$g_p\in H^0(X,H^{m-1})$ nonzero at $p$, taking $g_p=1$ when $m=1$.
The section $\ell_\eta g_p$ of $H^m$ vanishes at every support of $Z$.
Lemma~\ref{lem:annihilator} puts its germ in $\m_{X,p}^{s+1}H^m_p$.
The germ of $g_p$ is a unit, so division gives the displayed contact.
\end{proof}

This step explains why the factor $s^d$ in Theorem~\ref{thm:floor} is not an
extra geometric hypothesis.  It is forced by the combination of a common
tangent hyperplane and order-$s$ absorption.

\section{Milnor numbers and the dual multiplicity}

\begin{lemma}[Milnor floor from order]\label{lem:milnor}
Let $R=\mathbb C\{x_1,\ldots,x_d\}$ with maximal ideal $\m$, and let
$f\in\m^{s+1}$ have an isolated critical point at the origin.  Then
\[
 \mu(f)=\dim_{\mathbb C}R/
 \left(\frac{\partial f}{\partial x_1},\ldots,
 \frac{\partial f}{\partial x_d}\right)\geq s^d.
\]
Equality holds if the initial homogeneous form $f_{s+1}$ defines a smooth
hypersurface in $\PP^{d-1}$ (with the empty hypersurface convention for
$d=1$).
\end{lemma}

\begin{proof}
Let $J$ be the Jacobian ideal.  Isolatedness makes $J$ $\m$-primary.  It is
generated by a system of parameters in the regular local ring $R$, so
$\mu(f)=\operatorname{length}(R/J)=e(J)$.  Since $J\subseteq\m^s$,
monotonicity of Hilbert--Samuel multiplicity gives
\[
 e(J)\geq e(\m^s)=s^d.
\]
If $f_{s+1}$ is smooth projectively, its $d$ partial derivatives are a
homogeneous regular sequence of degree $s$.  The associated-graded
intersection has length $s^d$, so the local Jacobian colength is $s^d$.
\end{proof}

Dimca's multiplicity formula \cite{Dimca1986} says that for a smooth complex
hypersurface and a hyperplane section with isolated singularities,
\begin{equation}\label{eq:dimca}
 \mult_\eta(X^\vee)=\sum_{p\in\Sing(X\cap W_\eta)}\mu_p.
\end{equation}
The refined tangent-cone statement
\cite[Proposition~11.24]{Dimca1987} is
\begin{equation}\label{eq:tangent-cone}
 C_\eta(X^\vee)=\sum_p\mu_p[T_\eta(p^\perp)].
\end{equation}
Here $p^\perp$ is the hyperplane through $\eta$ in the dual projective space
corresponding, under biduality, to $p$; the displayed component is its linear
tangent space at $\eta$.  Equivalently, projectivizing gives the cycle
$\sum_p\mu_p[\PP(T_\eta(p^\perp))]$ in the projectivized ambient tangent
space.

For completeness, the branch contribution can also be read directly from
the Gauss map.  Choose analytic coordinates at $p$ in which $X$ is the graph
\[
 x_{d+1}=h(z_1,\ldots,z_d),\qquad h\in\m^{s+1},
\]
and $W_\eta=(x_{d+1}=0)$.  Work in the analytic local rings
$A=\cO^{\mathrm{an}}_{X^\vee,\eta}$ and
$R=\cO^{\mathrm{an}}_{X,p}\cong\mathbb C\{z_1,\ldots,z_d\}$, and put
$I=\m_\eta R$.  The local expression for the Gauss map gives, up to signs,
\begin{equation}\label{eq:gauss-ideal}
 I=\left(
 h_1,\ldots,h_d,\ h-\sum_{i=1}^d z_i h_i
 \right),\qquad h_i=\frac{\partial h}{\partial z_i}.
\end{equation}
Thus, for $J=(h_1,\ldots,h_d)$,
\[
 J\subseteq I\subseteq\m^s.
\]
When the initial form $h_{s+1}$ is projectively smooth, $J$ is a parameter
ideal with $e(J)=s^d$.  Monotonicity and $e(\m^s)=s^d$ then give
\[
 s^d=e(J)\geq e(I)\geq e(\m^s)=s^d,
\]
so the contribution of that normalization branch is exactly $s^d$.
The finite-extension associativity formula for Hilbert--Samuel multiplicity
recovers the sum over the points above $\eta$; see
\cite[Chapter~VIII, \S10, Corollary~1]{ZariskiSamuelII}.  Analytification and
completion preserve the multiplicities in question.  This is a local recovery
of the classical multiplicity--Milnor formula in this ordinary case, included
to make the equality mechanism in the examples explicit.

\begin{proof}[Proof of Theorem~\ref{thm:floor}]
Lemma~\ref{lem:normalization} gives the branch count.  Order-$s$ absorption
contains tangent absorption, so the exact point-span tangent floor
\cite[Theorem~1.1]{ErikssonTangent2026} gives
\[
 |Z|\geq N_{d,m}.
\]
Proposition~\ref{prop:contact} says that a local equation of the hyperplane
section belongs to $\m_{X,p}^{s+1}$.  The critical point is isolated, hence
Lemma~\ref{lem:milnor} gives $\mu_p\geq s^d$.  Equations
\eqref{eq:dimca} and \eqref{eq:tangent-cone} give all remaining assertions.

\end{proof}

\begin{remark}[Why the earlier rank term gives no refinement here]
Suppose $m\geq2s+1$.  Because $Z\subset W_\eta$, the equation of $W_\eta$
lies in the kernel of linear evaluation, so
\[
 r_1(Z):=\operatorname{rk}\bigl(H^0(X,H)\to H^0(Z,H|_Z)\bigr)\leq d+1.
\]
The quotient
\[
 \frac{\binom{d+2s+1}{d}}{\binom{d+s}{d}}
 =\prod_{j=1}^d\frac{2s+1+j}{s+j}
\]
is nondecreasing in $s\geq1$ term by term.  At $s=1$ it equals
$(d+3)(d+2)/6\geq d+1$.  Hence
\[
 \binom{d+m}{d}\geq(d+1)\binom{d+s}{d}
 \geq r_1(Z)\binom{d+s}{d}.
\]
Thus the rank-sensitive maximum in
\cite[Theorem~1.3]{ErikssonOsculating2026} specializes exactly to
$N_{d,m}$ for a Gauss fibre; presenting it as a stronger bound here would be
misleading.
\end{remark}

\section{Interpolation with prescribed ordinary jets}

We now adapt the proper-span construction of
\cite{ErikssonOsculating2026}; its antecedent for ordinary double points is
the author's public proof note \cite{ErikssonB219}.  The interpolation and
extension architecture $F=f+yG$ is therefore not claimed as new here.  The
added points are the prescribed ordinary initial forms, the exact Gauss-fibre
identity, and the resulting branch and dual-multiplicity calculation.

Let $W\cong\PP^d$ and
choose $N=N_{d,m}$ distinct points $Z\subset W$ such that
\begin{equation}\label{eq:unisolvent}
 H^0(W,\cO_W(m))\xrightarrow{\sim}H^0(Z,\cO_Z(m)).
\end{equation}
Such sets exist because the evaluation lines of all points span the dual of
$H^0(W,\cO_W(m))$; select a basis.  Very ampleness separates the supports.

\begin{lemma}[Simultaneous jets from separators]\label{lem:jets}
Let $p_1,\ldots,p_t$ be distinct points of $\PP^d$ and let $a_i\geq0$.
Restriction to the disjoint fat points is surjective in every degree
\[
 k\geq\sum_{i=1}^t(a_i+1)-1:
\]
\[
 H^0(\PP^d,\cO(k))\longrightarrow
 \bigoplus_{i=1}^t
 H^0\bigl(\cO(k)\otimes\cO_{\PP^d,p_i}/\m_{p_i}^{a_i+1}\bigr).
\]
\end{lemma}

\begin{proof}
For $i\ne j$, choose a linear form $\ell_{ij}$ vanishing at $p_j$ and
nonzero at $p_i$.  The product
\[
 P_i=\prod_{j\ne i}\ell_{ij}^{a_j+1}
\]
kills every other target and is a unit at $p_i$.  Degree-$a_i$ forms realize
all order-$a_i$ jets at $p_i$.  Their products with $P_i$ have common degree
$\sum_j(a_j+1)-1$ and isolate arbitrary data at $p_i$.  Summing over $i$
proves surjectivity in that degree.  Multiplication by a form nonzero at all
supports raises the degree.
\end{proof}

For every $p\in Z$, choose a homogeneous form
\[
 q_p\in\operatorname{Sym}^{s+1}(\m_{W,p}/\m_{W,p}^2)
\]
whose projective zero locus in $\PP(T_pW)$ is smooth.  Put
\[
 E=(s+2)N,
\]
and fix $D\geq E+1$.  Lemma~\ref{lem:jets} supplies
$f_0\in H^0(W,\cO_W(D))$ whose class at every $p$ is $q_p$ modulo
$\m_{W,p}^{s+2}$ and whose lower terms vanish.  Let
\[
 V=H^0(W,\cI_{(s+2)Z}(D)),\qquad \mathcal A=f_0+V.
\]

\begin{proposition}[A section with exactly the prescribed singularities]
\label{prop:section}
There is $f\in\mathcal A$ such that $V(f)\subset W$ is smooth away from $Z$.
At every $p\in Z$ it has an ordinary isolated hypersurface singularity of
multiplicity $s+1$ and Milnor number $s^d$.
\end{proposition}

\begin{proof}
Fix $q\in W\setminus Z$.  For every $p\in Z$, choose a hyperplane
$\ell_{p,q}$ through $p$ and not through $q$, and set
\[
 A_q=\prod_{p\in Z}\ell_{p,q}^{s+2}.
\]
It has degree $E$, belongs to the ideal of $(s+2)Z$, and is a unit on the
first neighbourhood $2q$.  Since $D-E\geq1$, multiplication by $A_q$ proves
that
\[
 V\longrightarrow H^0(2q,\cO_{2q}(D))
\]
is surjective.  Therefore the affine condition $j_q^1f=0$ has codimension
$d+1$ in $\mathcal A$.  Its incidence over the $d$-dimensional base
$W\setminus Z$ has dimension at most $\dim\mathcal A-1$.  Choose $f$ outside
the closure of its image.  The fixed jets $q_p$ give the asserted ordinary
singularities, and Lemma~\ref{lem:milnor} gives $\mu_p=s^d$.
\end{proof}

\section{Smooth hypersurfaces and exact Gauss fibres}

Embed $W=V(y)$ as a hyperplane in $\PP^{d+1}$ and lift $f$ independently of
$y$.  Put
\[
 U=H^0(\PP^{d+1},\cO(D-1)),\qquad F_G=f+yG\quad(G\in U).
\]

\begin{proposition}[Smooth extension]\label{prop:extension}
There is $G\in U$, nonzero at every point of $Z$, such that
$X=V(F_G)$ is a smooth integral hypersurface.  Moreover
\[
 \Sing(X\cap W)=Z,
 \qquad
 (\gamma_X^{-1}([W]))_{\mathrm{red}}=Z.
\]
\end{proposition}

\begin{proof}
On $y\ne0$, multiplication by $y$ is a unit on first neighbourhoods and
$\cO(D-1)$ generates first jets.  For fixed $q$ in that $(d+1)$-dimensional
open set, the affine condition $j_q^1(f+yG)=0$ has codimension $d+2$ in $U$.
The corresponding incidence has dimension at most $\dim U-1$, so the closure
of its image is proper.  Avoid it together with the finitely many hyperplanes
$G(p)=0$.

The resulting $X$ is smooth on $y\ne0$.  At a point of $W\setminus Z$ lying
on $X$, the tangential first jet of $f$ is nonzero by
Proposition~\ref{prop:section}.  At $p\in Z$, the normal first jet is
$G(p)\ne0$.  Thus $X$ is smooth everywhere.  The hypersurface sequence and
$H^1(\PP^{d+1},\cO(-D))=0$ give $H^0(X,\cO_X)=\mathbb C$, so $X$ is
connected; a smooth connected scheme is integral.

The hyperplane section is $V(f)$, hence has singular locus exactly $Z$.
A point of $X\cap W$ belongs to the Gauss fibre over $[W]$ precisely when the
section is singular there.  A point outside $W$ cannot have $W$ as tangent
hyperplane because a tangent hyperplane contains its point.  This proves the
exact fibre identity.
\end{proof}

At $p\in Z$, the local equation on $X$ is
\[
 y=-f/G\in\m_{X,p}^{s+1},
\]
and its initial form is the nonzero scalar multiple $-q_p/G(p)$.  Thus
$X\cap W$ has precisely the ordinary singularities already prescribed.

\begin{proof}[Proof of Theorem~\ref{thm:sharpness}]
Use Proposition~\ref{prop:extension}.  Since $D>m$, restriction is an
isomorphism
\[
 H^0(\PP^{d+1},\cO(m))\xrightarrow{\sim}H^0(X,\cO_X(m)).
\]
Let a section on $X$ vanish on $Z$ and let $Q$ be its ambient lift.  By
\eqref{eq:unisolvent}, $Q|_W=0$, so $Q=yR$.  Since
$y\in\m_{X,p}^{s+1}$ at every support, $Q$ vanishes on $(s+1)Z$.  The reverse
kernel inclusion is automatic, hence the annihilator criterion proves
order-$s$ absorption.  Unisolvence gives
\[
 |Z|=\dim S_Z=N.
\]
The ambient degree-$m$ space has dimension $\binom{d+m+1}{d+1}>N$, so the
point span is proper.

The exact Gauss fibre is Proposition~\ref{prop:extension}; the section
singularities have Milnor number $s^d$.  Dimca's formulas
\eqref{eq:dimca}--\eqref{eq:tangent-cone} give the stated dual multiplicity
and tangent-cone cycle.  When $s=1$, each section singularity is $A_1$, the
Gauss map is immersive at its support, and the corresponding dual branch is
smooth.  The tangent hyperplanes $p^\perp$ are distinct because the supports
are distinct.
\end{proof}

\section{Context, reproducibility, and limitations}

The identity $\mult_W(X^\vee)=\sum\mu_p$ and its tangent-cone refinement are
not new; they are attributed respectively to Dimca
\cite{Dimca1986} and \cite[Proposition~11.24]{Dimca1987}, with broader
multiplicity results due to Parusi\'nski \cite{Parusinski1991}, and belong to
the classical projective-duality and discriminant framework
\cite{Kleiman1986,GKZ1994}.  Dimca--Ilardi prove that a generic smooth
hypersurface in $\PP^n$, $n\geq3$, has a dual normal-crossing point of
multiplicity $n$, arising from $n$ nodes in a hyperplane section
\cite{DimcaIlardi2024}.  They also emphasize the classical fact that isolated
singularity hyperplane sections can be extended to smooth hypersurfaces.
The present construction uses that freedom subject to the additional
unisolvence and absorption constraints.

The normalizing role of the Gauss map and the earlier large-fibre construction
also occur in the author's public notes \cite{ErikssonB218,ErikssonB219}.
Those notes do not state the branch count or dual multiplicity computed here;
they are nevertheless direct authorial antecedents rather than independent
validation.  The exact tangent floor is invoked from
\cite{ErikssonTangent2026}; the higher-block estimate and the underlying
proper-span construction are invoked from \cite{ErikssonOsculating2026}, to
which the exact reduced Gauss-fibre and dual-singularity calculations are
added here.  These are papers by the same author and are not independent
external validation.

The accompanying replay
\texttt{repro/verify\_dual\_multiplicity.py} checks finite monomial Jacobian
quotients, 105 binomial/multiplicity cases, 100 rank-collapse cases, and
selected simplex evaluation matrices using integer and rational arithmetic.
The byte-stable runner is
\texttt{repro/run\_all\_replays.py}.  These scripts do not mechanize the
incidence arguments, projective duality, Dimca's formula, or literature
priority.

The limitations are material.
\begin{itemize}
\item The dual-singularity theorem is over $\mathbb C$.  No positive-
characteristic Milnor, inseparable Gauss-map, or wild-ramification extension
is asserted.
\item The support is the complete reduced Gauss fibre, not an arbitrary
subset.  Its scheme-theoretic fibre can be nonreduced and is not classified.
\item The degree $D\geq(s+2)N_{d,m}+1$ is a transparent sufficient threshold,
not a minimality result.
\item For $s=1$ the tangent cone is a reduced arrangement of distinct
hyperplanes.  It is not called a simple normal-crossing divisor when the
number of branches exceeds the ambient dimension.
\item No analytic classification of the dual branches, equality
classification, exhaustive priority review, human peer review, or formal
verification is claimed.
\end{itemize}

\section{Conclusion}

Osculating absorption in a Gauss fibre forces more than many coincident
tangent points.  It forces the common tangent hyperplane to have order
$s+1$ contact at every support, converting the exact point-span rank floor
into the sharp dual-singularity floor
\[
 \mult_W(X^\vee)\geq s^d\binom{d+m}{d}.
\]
The previously developed prescribed-jet construction realizes equality while
keeping the point span proper and the reduced Gauss fibre exact; the new use
made of it here is the dual-singularity calculation.  The result ties finite
reduced span rank,
osculating contact, Milnor numbers, normalization branches, and dual
multiplicity in one sharp construction.  Minimal ambient degree and the
geometry of equality configurations remain open.

{\small
\begin{thebibliography}{99}

\bibitem{Dimca1986}
A.~Dimca,
\emph{Milnor numbers and multiplicities of dual varieties},
Rev. Roumaine Math. Pures Appl. \textbf{31} (1986), no.~6, 535--538.

\bibitem{Dimca1987}
A.~Dimca,
\emph{Topics on Real and Complex Singularities},
Advanced Lectures in Mathematics, Friedr. Vieweg \& Sohn, Braunschweig, 1987.

\bibitem{DimcaIlardi2024}
A.~Dimca and G.~Ilardi,
\emph{On the duals of smooth projective complex hypersurfaces},
Publ. Mat. \textbf{68} (2024), 431--438.
\url{https://doi.org/10.5565/PUBLMAT6822404}

\bibitem{Parusinski1991}
A.~Parusi\'nski,
\emph{Multiplicity of the dual variety},
Bull. London Math. Soc. \textbf{23} (1991), no.~5, 429--436.
\url{https://doi.org/10.1112/blms/23.5.429}

\bibitem{ErikssonTangent2026}
L.~Eriksson,
\emph{The exact rank floor for point-span tangent absorption in arbitrary
characteristic}, ARR-2026-2MHNZRRJP49Y9SWP, v3 (2026).
\url{https://arr-research.github.io/papers/ARR-2026-2MHNZRRJP49Y9SWP/versions/v3/}

\bibitem{ErikssonOsculating2026}
L.~Eriksson,
\emph{Exact floors and proper-span extremizers for higher osculating
absorption}, ARR-2026-66Q8M61AA196T8BC, v2 (2026).
\url{https://arr-research.github.io/papers/ARR-2026-66Q8M61AA196T8BC/versions/v2/}

\bibitem{ErikssonB218}
L.~Eriksson,
\emph{The ordinary Gauss map is finite birational}, proof note B218 (2026).
\url{https://github.com/lluiseriksson/hodge-conjecture-research-front/blob/main/proofs/B218-gauss-map-finite-birational.md}

\bibitem{ErikssonB219}
L.~Eriksson,
\emph{Special Gauss fibers can be arbitrarily large}, proof note B219 (2026).
\url{https://github.com/lluiseriksson/hodge-conjecture-research-front/blob/main/proofs/B219-arbitrarily-large-special-gauss-fibers.md}

\bibitem{GKZ1994}
I.~M.~Gelfand, M.~M.~Kapranov, and A.~V.~Zelevinsky,
\emph{Discriminants, Resultants, and Multidimensional Determinants},
Birkh\"auser, Boston, 1994.

\bibitem{Kleiman1986}
S.~L.~Kleiman,
\emph{Tangency and duality}, in
\emph{Proceedings of the 1984 Vancouver Conference in Algebraic Geometry},
CMS Conf. Proc. \textbf{6}, Amer. Math. Soc., 1986, 163--225.

\bibitem{ZariskiSamuelII}
O.~Zariski and P.~Samuel,
\emph{Commutative Algebra, Volume II},
Graduate Texts in Mathematics 29, Springer, 1975.
\url{https://doi.org/10.1007/978-3-662-29244-0}

\end{thebibliography}
}

\end{document}
